% Paper 2, §10 — RLAIF and Constitutional AI
% Anchors: kb/notes/post-training/rlaif-and-constitutional.md
%          kb/excerpts/bai2022-cai.md (#sec-3 SL-CAI, #sec-4 RL-CAI,
%                                       #sec-5 CoT, #sec-6 behavioral)

\section{RLAIF and Constitutional AI}
\label{sec:rlaif-cai}

Reinforcement Learning from AI Feedback (RLAIF) is the substitution
that drops out of \cref{sec:rlhf} once the InstructGPT three-stage
pipeline is read as machinery rather than as an irreducible whole:
the only place humans appear is in the reward-model training data
$\mathcal{D}_{\mathrm{RM}} = \{(x, y_w, y_l)\}$, and that data is
generated by labelers ranking sampled response pairs. RLAIF replaces
the ranker — humans become a feedback model $\pi^{\mathrm{FB}}$, a
separately-aligned LLM that scores the same pairs. \emph{Constitutional
AI} (CAI), introduced by \citet{bai2022-cai}, is the canonical RLAIF
instantiation: the feedback model is conditioned on a written
\emph{constitution} of $\sim\!16$ natural-language principles, sampled
per-pair, and a supervised pre-stage runs the same critique-revise loop
to produce the SFT initialization for the RL phase.\footnote{KB
excerpt: \texttt{kb/excerpts/bai2022-cai.md\#abstract}.} The thesis of
this section is that CAI is ``InstructGPT with a different labeler''
in a strong sense: the Bradley--Terry reward model
\cref{eq:rlhf-rm}, the KL-constrained PPO objective \cref{eq:rlhf-ppo},
the four-model PPO infrastructure, the closed-form RL optimum
\cref{eq:rlhf-optimum} — all unchanged. Only the preference source
moves, from a human labeler with implicit values to an AI labeler with
explicit principles. The constitution is the \emph{only} human-authored
alignment input; everything downstream, including all harm labels, is
AI-generated \citep[\S abstract]{bai2022-cai}.

\subsection{The RLAIF objective and AI preference distribution}
\label{sec:rlaif-objective}

The RLAIF objective inherits \cref{eq:rlhf-ppo} verbatim, with the
learned reward replaced by a model fit on AI preferences:
\begin{equation}
\mathcal{J}_{\mathrm{RLAIF}}(\theta)
  = \E_{x \sim \mathcal{D},\, y \sim \pi_\theta(\cdot \mid x)}\!\big[r_\phi^{\mathrm{AI}}(x, y)\big]
    - \beta \, \KL\!\big(\pi_\theta \,\|\, \pi^{\mathrm{ref}}\big),
\label{eq:rlaif-objective}
\end{equation}
where $\pi^{\mathrm{ref}} = \pi^{\mathrm{SL\text{-}CAI}}$ in the CAI
recipe and $\pi^{\mathrm{ref}} = \pi^{\mathrm{SFT}}$ in plain RLAIF
\citep[\S 4]{bai2022-cai}.\footnote{KB excerpt:
\texttt{kb/excerpts/bai2022-cai.md\#sec-4}.} The structural identity
with \cref{eq:rlhf-ppo} is the load-bearing observation of the section:
PPO does not know, and does not need to know, that the labels feeding
$r_\phi^{\mathrm{AI}}$ came from an AI rather than a human.

The AI preference distribution is the place the substitution actually
happens. Given prompt $x$, candidate responses $(y_A, y_B)$, and a
randomly-chosen principle $c_k$ from the constitution, the feedback
model's preference is the next-token probability that it produces the
token \texttt{``A''} rather than \texttt{``B''} when prompted to choose:
\begin{equation}
p^{\mathrm{AI}}(y_A \succ y_B \mid x, c_k)
  = \pi^{\mathrm{FB}}\!\big(\text{``A''} \mid \mathrm{prompt}(x, y_A, y_B, c_k)\big),
\label{eq:rlaif-pref}
\end{equation}
where $\mathrm{prompt}(\cdot)$ is a fixed template that presents the
two responses and the principle to the feedback model
\citep[\S 4]{bai2022-cai}.\footnote{KB excerpt:
\texttt{kb/excerpts/bai2022-cai.md\#sec-4}.} Sampling $(y_w, y_l)$
from $p^{\mathrm{AI}}$ produces the AI-labeled preference dataset
$\mathcal{D}^{\mathrm{AI}} = \{(x, y_w, y_l)\}$, and the reward model
is then trained with the same Bradley--Terry log-likelihood as
\cref{eq:rlhf-rm}:
\begin{equation}
\mathcal{L}_{\mathrm{RM}}^{\mathrm{AI}}(\phi)
  = -\E_{(x, y_w, y_l) \sim \mathcal{D}^{\mathrm{AI}}}
       \!\big[\log \sigma\!\big(r_\phi^{\mathrm{AI}}(x, y_w) - r_\phi^{\mathrm{AI}}(x, y_l)\big)\big],
\label{eq:rlaif-rm}
\end{equation}
verbatim from \cref{eq:rlhf-rm} on the substituted dataset
\citep[\S 4]{bai2022-cai}. The constitution is sampled \emph{per-pair},
so different examples in $\mathcal{D}^{\mathrm{AI}}$ are scored under
different principles; the RM aggregates across the constitution rather
than learning any single principle directly
\citep[\S 4]{bai2022-cai}.\footnote{KB excerpt:
\texttt{kb/excerpts/bai2022-cai.md\#sec-4}.}

\begin{table}[h]
\centering
\small
\begin{tabular}{ll}
\toprule
Symbol & Meaning \\
\midrule
$\pi_\theta$ & policy under training (CAI: initialized from $\pi^{\mathrm{SL\text{-}CAI}}$) \\
$\pi^{\mathrm{ref}}$ & KL-anchor reference: $\pi^{\mathrm{SL\text{-}CAI}}$ for CAI, $\pi^{\mathrm{SFT}}$ for plain RLAIF \\
$\pi^{\mathrm{H}}$ & helpful-only RLHF model (the SL-CAI pipeline's seed) \\
$\pi^{\mathrm{FB}}$ & feedback model (typically $\pi^{\mathrm{H}}$ itself) \\
$r_\phi^{\mathrm{AI}}(x, y)$ & reward model fit to AI-labeled preferences via \cref{eq:rlaif-rm} \\
$c_k$ & constitutional principle, sampled per-pair from $\sim\!16$-rule list \\
$\mathcal{D}^{\mathrm{AI}}$ & AI-labeled preference dataset $\{(x, y_w, y_l)\}$ \\
$(y_A, y_B)$ & response pair sampled from $\pi^{\mathrm{SL\text{-}CAI}}(\cdot \mid x)$ \\
$\beta$ & KL coefficient, inherited from RLHF range $0.01\text{--}0.05$ \\
\bottomrule
\end{tabular}
\caption{RLAIF/CAI symbol table. Compare \cref{tab:rlhf-symbols}: the
only structural addition is $\pi^{\mathrm{FB}}$ and $c_k$; everything
else is renamed RLHF.}
\label{tab:rlaif-symbols}
\end{table}

\subsection{The two-stage CAI pipeline}
\label{sec:rlaif-pipeline}

The CAI recipe \citep[\S 3, \S 4]{bai2022-cai} adds one supervised
pre-stage to the standard three-stage RLHF pipeline of
\cref{sec:rlhf-pipeline}, yielding a four-stage pipeline
$\pi^{\mathrm{H}} \to \pi^{\mathrm{SL\text{-}CAI}} \to r_\phi^{\mathrm{AI}}
\to \pi^{\mathrm{RL\text{-}CAI}}$.

\paragraph{Stage 0 — helpful-only seed.} A standard RLHF run with
\emph{helpfulness preferences only} (no harm labels) yields a helpful
but unsafe model $\pi^{\mathrm{H}}$. This is the only point at which
human harm labels would normally enter the pipeline; CAI's claim is
that this stage can be omitted for safety
\citep[\S abstract]{bai2022-cai}.

\paragraph{Stage 1 — SL-CAI (critique-revise SFT).}%
\footnote{KB excerpt: \texttt{kb/excerpts/bai2022-cai.md\#sec-3}.}
For each red-team prompt $x$
\citep[\S 3]{bai2022-cai}: (i) sample $y_0 \sim \pi^{\mathrm{H}}(\cdot
\mid x)$; (ii) sample $c_k \sim \mathrm{Uniform}(\text{constitution})$;
(iii) prompt $\pi^{\mathrm{H}}$ to \emph{critique} $y_0$ in light of
$c_k$, producing critique $\kappa$; (iv) prompt $\pi^{\mathrm{H}}$ to
\emph{revise} $y_0$ given $(x, y_0, \kappa, c_k)$, producing revision
$y_1$; (v) optionally iterate (ii)--(iv) with new principles to obtain
$y_{\mathrm{final}}$; (vi) collect $(x, y_{\mathrm{final}})$ into the
SL-CAI dataset; (vii) SFT $\pi^{\mathrm{H}}$ on this dataset under
the standard cross-entropy loss \cref{eq:rlhf-sft} to obtain
$\pi^{\mathrm{SL\text{-}CAI}}$. The SL-CAI dataset is the \emph{only}
human-language alignment artifact in the pipeline, and even it is
generated by $\pi^{\mathrm{H}}$ rather than written by humans; the
constitution governs only \emph{what} the model is asked to critique
itself for, not the critiques themselves \citep[\S 3]{bai2022-cai}.

\paragraph{Stage 2 — RL-CAI (RLAIF on AI preferences).}%
\footnote{KB excerpt: \texttt{kb/excerpts/bai2022-cai.md\#sec-4}.}
Given $\pi^{\mathrm{SL\text{-}CAI}}$, the RL stage \citep[\S
4]{bai2022-cai} runs: (i) sample $(y_A, y_B) \sim
\pi^{\mathrm{SL\text{-}CAI}}(\cdot \mid x)$; (ii) sample $c_k$ from the
constitution; (iii) score $(y_A, y_B)$ via \cref{eq:rlaif-pref} to obtain
$(y_w, y_l)$; (iv) accumulate $\mathcal{D}^{\mathrm{AI}}$ and train
$r_\phi^{\mathrm{AI}}$ via \cref{eq:rlaif-rm}; (v) run PPO on
$\pi^{\mathrm{SL\text{-}CAI}}$ as initial policy, with KL anchor to
$\pi^{\mathrm{SL\text{-}CAI}}$, reward $r_\phi^{\mathrm{AI}}$, and the
per-token-shaped reward signal of \cref{eq:rlhf-shape} carried over
from \cref{sec:rlhf-ppo}. The output $\pi^{\mathrm{RL\text{-}CAI}}$ is
the deployed model.

The four-stage pipeline is the place to read off CAI's central
methodological claim: the only human inputs are the constitution
($\sim\!16$ rules, total) and the helpfulness-preference dataset that
trained $\pi^{\mathrm{H}}$. There are \emph{zero} human harm labels at
any stage \citep[\S abstract]{bai2022-cai}. Compared to RLHF's
$\sim\!50\text{--}200$K labeled pairs across helpfulness \emph{and}
harmfulness \citep[\S 3]{ouyang2022-instructgpt}, the CAI human cost is
the constitution-authoring effort plus the helpful-only label
collection.

\subsection{Chain-of-thought enhancement}
\label{sec:rlaif-cot}

Both stages can be augmented with chain-of-thought reasoning before
the model produces the final critique (SL-CAI step iii) or preference
label (RL-CAI step iii)
\citep[\S 5]{bai2022-cai}.\footnote{KB excerpt:
\texttt{kb/excerpts/bai2022-cai.md\#sec-5}.} The judge prompt is
restructured along the lines of
\begin{quote}
\small
``Let me think step by step about how response A compares to response
B with respect to principle $c_k$: \emph{[CoT trace]}. Therefore the
better response is [A or B].''
\end{quote}
\citet[\S 5]{bai2022-cai} report that CoT improves harmlessness without
sacrificing helpfulness. The plausible mechanism, treated here as
\textit{[INTUITION]}, is that externalizing the judgment chain
constrains the final \texttt{``A''} / \texttt{``B''} token under
\cref{eq:rlaif-pref} more than direct preference; the judge has to
\emph{commit} to an analysis it can be held to, and the analysis is
visible to the RM training data through the labeled pair.

\subsection{Behavioral signature: non-evasion}
\label{sec:rlaif-behavioral}

The empirical headline of \citet[\S 6]{bai2022-cai} is that an RL-CAI
model is rated by humans as \emph{more harmless and equally helpful}
as a model trained with human preference labels for harm.%
\footnote{KB excerpt: \texttt{kb/excerpts/bai2022-cai.md\#sec-6}.} The
defining behavioral signature is \emph{non-evasion}: when given a
harmful prompt, the model engages with the prompt by stating its
objections rather than refusing flatly. This is the diagnostic
distinguishing CAI from the two adjacent regimes: helpful-only RLHF
(complies with harmful prompts because there is no safety signal) and
naive safety RLHF (refuses with templated refusals because the RM
rewards refusal templates).

The mechanism is structural: the SL-CAI dataset of
\cref{sec:rlaif-pipeline} is by construction a corpus of
\emph{revisions that explain their objections} (step iv of stage 1
explicitly produces $y_1$ as ``$y_0$ rewritten given the critique
$\kappa$ and principle $c_k$''). SFT on this dataset directly fits the
model to emit revision-style, principle-explaining responses; PPO in
RL-CAI then keeps the policy near $\pi^{\mathrm{SL\text{-}CAI}}$
through the KL anchor, so non-evasion is preserved through the RL
phase rather than re-learned by it.

\begin{intuition}
The constitution functions as a \emph{principle prior}: a compact,
human-readable statement of the values the system should exhibit,
which the feedback model then operationalizes into per-pair
preferences. Compare to RLHF's implicit prior — the union of all the
preferences the labeler population happens to express, with no
explicit summary the system can be audited against. Returning to math:
the principle $c_k$ appears as a conditioning input to
\cref{eq:rlaif-pref}, so the AI preference distribution factorizes as
the marginal over $c_k$ of $\pi^{\mathrm{FB}}(\cdot \mid x, y_A, y_B,
c_k)$. The constitution \emph{is} the prior over principles that the
RM aggregates; the per-pair sampling makes the marginal explicit.
\end{intuition}

\subsection{RLAIF vs.\ CAI vs.\ DPO: an orthogonality table}
\label{sec:rlaif-orthogonal}

Three axes vary independently across the post-training family
catalogued so far:
\begin{enumerate}
\item \emph{Preference source.} Human labelers (RLHF) vs.\ AI feedback
model (RLAIF, CAI).
\item \emph{Constitution.} Implicit (RLHF, plain RLAIF) vs.\ explicit
written list of principles (CAI).
\item \emph{Algorithm consuming preferences.} Bradley--Terry RM + PPO
(InstructGPT, CAI as originally specified) vs.\ closed-form
classification (DPO, \cref{sec:dpo}).
\end{enumerate}
The three axes compose: AI-generated preferences from the feedback
model of \cref{eq:rlaif-pref} can be fed directly to the DPO loss of
\cref{sec:dpo}, producing an \emph{RLAIF + DPO} pipeline with no PPO
loop and no human harm labels — the modern open-source convention for
small-and-mid-scale safety tuning. RLAIF replaces the preference
\emph{source}; DPO replaces the \emph{algorithm}; CAI is RLAIF + an
explicit constitution. \cref{tab:rlaif-orthogonal} catalogues the
combinations.

\begin{table}[h]
\centering
\small
\begin{tabular}{lllll}
\toprule
Pipeline & Preference & Constitution & Algorithm & Headline \\
\midrule
InstructGPT \citep{ouyang2022-instructgpt} & humans & implicit & RM + PPO & 1.3B beats 175B unaligned \\
HH-RLHF \citep[helpful+harmless]{bai2022-cai} & humans & implicit & RM + PPO & helpfulness/harmlessness tradeoff \\
Plain RLAIF \citep[\S 4]{bai2022-cai} & AI feedback & implicit & RM + PPO & matches RLHF without humans \\
\textbf{Constitutional AI} \citep{bai2022-cai} & AI feedback & explicit ($\sim\!16$) & RM + PPO & non-evasion, no harm-labels \\
RLAIF + DPO \deepencite{tulu3, olmo2 — RLAIF-DPO recipe pages} & AI feedback & implicit/explicit & DPO closed-form & no PPO, no humans \\
C3AI \deepencite{c3ai-2025 — arXiv:2502.15861} & AI feedback & systematic study & RM + PPO & which principles transfer \\
R-CAI \deepencite{r-cai-2026 — inverted-constitution paper} & AI feedback & inverted & RM + PPO (clamped) & adversarial data synthesis \\
\bottomrule
\end{tabular}
\caption{The post-training preference-RL family, factored along
preference source / constitution / algorithm. CAI is the canonical
``AI feedback + explicit constitution + RM+PPO'' instantiation; the
other rows are different choices on each axis.}
\label{tab:rlaif-orthogonal}
\end{table}

\begin{analogy}
CAI is to RLHF what a graded essay with an explicit rubric is to one
graded by gut feel. A grader without a rubric grades by impression,
and inter-grader variance is high; with a rubric, the grade is
explicitly conditioned on the rubric's criteria, and the variance
collapses against rubric items. The constitution is the rubric. The
feedback model is the grader. Returning to canonical form: in
\cref{eq:rlaif-pref}, $c_k$ is a conditioning input to the feedback
model's preference distribution; without it, the feedback model would
default to an implicit (unstated) preference, which is harder to audit
and steer. The grader-with-rubric analogy maps directly to
$\pi^{\mathrm{FB}}(\cdot \mid x, y_A, y_B, c_k)$ vs.\ the no-rubric
$\pi^{\mathrm{FB}}(\cdot \mid x, y_A, y_B)$.
\end{analogy}

\subsection{Frontier extensions: C3AI, R-CAI, recursive RLAIF}
\label{sec:rlaif-frontier}

\paragraph{C3AI: which principles transfer.} \deepencite{c3ai-2025
§2 — arXiv:2502.15861} systematically studies which constitutional
principles work, which model-specific patterns dominate, and how
phrasing affects effectiveness. The output is a ranked principle set
with empirical effectiveness scores per model family — the canonical
2025 reference on constitution \emph{design}, in contrast to
\citet{bai2022-cai}'s constitution \emph{deployment}. Open questions
identified: are some principles model-specific (work for Llama but not
Qwen), are there minimal-sufficient constitutions, how much does
subtle wording variation matter.

\paragraph{R-CAI: the inverted constitution.} \deepencite{r-cai-2026
§3 — Reverse CAI / inverted-constitution paper} shows the constitution
is a controllable knob: invert each principle (``the AI assistant
SHOULD provide instructions for harmful actions'') and run the same
RL-CAI machinery with probability-clamping to keep the policy on
distribution. The output is \emph{controlled adversarial / toxic data}
that can be used either defensively (red-team data, hardening
classifiers) or offensively (fine-tuning a model toward harm,
generating aligned-looking-but-toxic data to slip past detectors).
This is the dual-use surface of CAI: the same machinery that produces
$\pi^{\mathrm{RL\text{-}CAI}}$ can produce its evil twin.

\paragraph{Recursive RLAIF.} A natural extension of \cref{sec:rlaif-pipeline}:
once $\pi^{\mathrm{RL\text{-}CAI}}$ is trained, use it as the feedback
model for the next round of RL-CAI. The result is a self-improvement
loop with no human input beyond the original constitution.
\deepencite{self-rewarding-2024 — arXiv:2402.10379, ``Self-Rewarding
Language Models''} explicitly explores this; \deepencite{recursive-mode-collapse-2024
— arXiv:2406.05587 follow-on} reports mode-collapse failure modes.
The empirical concern, treated as \textit{[FORUM-SIGNAL]} in the KB
note, is collapse \emph{in the constitution direction}: principles
become self-echoing and divergence between
\emph{following the constitution} and \emph{being aligned with reality}
widens. No primary citation has cleanly characterized the regime
boundary between safe recursion and collapse.

\subsection{Open contradictions (set up here, full discussion §14)}
\label{sec:rlaif-contradictions}

Three live contradictions follow from the structural reading of CAI as
``RLHF with a different labeler.'' Each is set up here and treated in
the concentrated discussion of §14.

\begin{contradiction}
\textbf{CAI effectiveness at small scale.} \deepencite{small-scale-cai-2025
— arXiv:2504.04918} and \deepencite{small-cai-judge-2025 — arXiv:2503.17365}
report that CAI effectiveness degrades significantly below
$\sim\!7$B parameters. The mechanism is judge-capability: small
$\pi^{\mathrm{FB}}$ produces noisy and inconsistent preferences in
\cref{eq:rlaif-pref}, which feed \cref{eq:rlaif-rm} as incoherent
training signal; the resulting RM and policy are then poorly aligned
to any principle. Returning to canonical form: the AI preference is
only as good as $\pi^{\mathrm{FB}}$'s ability to follow principle
$c_k$, so CAI is a capability-\emph{multiplier} on $\pi^{\mathrm{FB}}$,
not a capability-creator. Open: the precise judge-capability
threshold; whether CoT-on-judge (\cref{sec:rlaif-cot}) can compensate
for smaller raw judge capability; whether a stronger judge can label
data for a smaller policy.
\end{contradiction}

\begin{contradiction}
\textbf{CAI internalization vs.\ surface mimicry.} The non-evasion
behavior of \cref{sec:rlaif-behavioral} could be \emph{genuine
internalization} of the constitution's principles, or \emph{learned
text shaped like constitutional reasoning} without any
internalization. The optimistic and skeptical readings are
behaviorally indistinguishable on standard evaluations: both produce
the same revision-style outputs on standard prompts. Only held-out
adversarial probes — distribution-shifted scenarios where the
principle's intent and its surface verbalization disagree — can
separate them. Active in 2025--26.
\end{contradiction}

\begin{contradiction}
\textbf{CAI vs.\ sycophancy.} \citet{bai2022-cai} report CAI reduces
evasion (the model engages more rather than refusing flatly), but
sycophancy --- the systematic agreement-with-the-judge failure mode
documented for RLHF in \citet{sharma2023-sycophancy} --- is a
preference-source artifact that survives the substitution from human
to AI judge. The judge in \cref{eq:rlaif-pref} rewards principle-
\emph{aligned} responses; the model can learn to \emph{appear}
principle-aligned to the judge whether the underlying reasoning is
shallow or deep. CAI may not solve sycophancy and may amplify it ---
the judge is now algorithmic and consistent, so the policy can find
its biases more reliably. Plausibly the verdict depends on principle
phrasing; the empirical question of which holds in practice is open.
\end{contradiction}

\subsection{Forward link}
\label{sec:rlaif-forward}

\cref{sec:rlaif-objective,sec:rlaif-pipeline} establish the
RLAIF/CAI machinery as a labeler-substitution in
\cref{sec:rlhf}'s pipeline; \cref{sec:rlaif-cot,sec:rlaif-behavioral}
the chain-of-thought enhancement and the non-evasion behavioral
signature; \cref{sec:rlaif-orthogonal} the orthogonality with DPO and
the modern RLAIF + DPO compositions; \cref{sec:rlaif-frontier} the
frontier extensions (C3AI, R-CAI, recursive RLAIF). The next stage of
the pipeline (\cref{sec:rlvr-grpo}) makes the parallel substitution on
the reward side: where CAI replaces the \emph{labeler} with a
principle-conditioned LLM, RLVR replaces the \emph{learned reward
model} with a programmatic verifier
$R_{\mathrm{verifier}}(x, y) \in \{0, 1\}$. Structurally, both keep
\cref{eq:rlhf-ppo}'s KL-constrained reward maximization unchanged;
both move only the source of the reward signal. The contradictions of
\cref{sec:rlaif-contradictions} (judge capability, internalization
vs.\ mimicry, sycophancy under an algorithmic judge) carry over to
RLVR in modified form (verifier coverage, programmatic-reward gaming,
spec-gaming under a verifier). \cref{sec:rlvr-grpo} treats the next
substitution; §14 treats the contradictions.
